\label{sec:observables}

Among the simplest quantities one can consider within the \gff\ are
vacuum expectation values of gauge-invariant operators at finite flow
time. It is one of the remarkable properties of the \gff\ that these
operators do not require any renormalization beyond that of regular
\qcd, and that of the involved flowed fields. For the gauge coupling and
the quark masses, the \msbar\ renormalization is given by the
replacement
\begin{equation}
  \begin{split}
    g\to g_0\equiv \left(\frac{\mu\mathrm{e}^{\gamma_\mathrm{E}/2}}{\sqrt{4\pi}}\right)^\ep
    Z_g(\alphas(\mu))\,g(\mu)\,,\qquad m_f\to m_{f,0}\equiv
    Z_m(\alphas(\mu))\,m_f(\mu)
    \label{eq:renorm}
  \end{split}
\end{equation}
in \eqn{eq:lags}, where $\mu$ is the renormalization scale,
$\alphas=g^2/(4\pi)$, and $\EulerGamma=0.5772\ldots$ the
Euler-Mascheroni constant.  Through the perturbative order required in
this paper, the renormalization constants are given by
\begin{equation}
  \begin{split}
    Z_g(\alphas) &=
    1 -
    \frac{\alphas}{4\pi}\frac{\beta_0}{2\ep} +
    \left(\frac{\alphas}{4\pi}\right)^2\left(\frac{3\beta_0^2}{8\ep^2}
    -\frac{\beta_1}{4\ep}\right) + \order{\alphas^3}\,, \\ Z_m(\alphas) &=
    1 -\frac{\alphas}{4\pi}\frac{\gamma_{m,0}}{2\ep}
    +\left(\frac{\alphas}{4\pi}\right)^2\left[\frac{1}{\ep^2}\left(
      \frac{\gamma_{m,0}^2}{8} + \frac{\beta_0\gamma_{m,0}}{4}\right)
      -\frac{\gamma_{m,1}}{4\ep}\right] + \order{\alphas^3}\,,
    \label{eq:rgconsts}
  \end{split}
\end{equation}
with
\begin{equation}
\begin{aligned}
      \beta_0 &= \frac{11}{3} \CA - \frac{4}{3} \TF\,, &
      \beta_1 &= \frac{34}{3} \CA^2 - \left( 4 \CF + \frac{20}{3} \CA
      \right) \TF \,,\\
      \gamma_{m,0} &= 6 \CF\,, &
      \gamma_{m,1} &= \frac{97}{3} \CA \CF + 3 \CF^2
      - \frac{20}{3} \CF \TF \,.
    \label{eq:anombg}
\end{aligned}
\end{equation}
$\CF$ and $\CA$ are the quadratic Casimir eigenvalues of the fundamental
and the adjoint representation of the gauge group,
respectively. Furthermore, $\TF=\TR \NF$, with $\NF$ the number of quark
flavors, and $\TR$ the trace normalization in the fundamental
representation. For SU($\NC$), it is $\CF=(\NC^2-1)/(2\NC)$, $\CA=\NC$,
and $\TR=1/2$.

The flowed gauge field $B^a_\mu(t,x)$ does not require renormalization,
so that for example matrix elements of the gluon action density,
\begin{align}
  E (t,x) \equiv \frac{1}{4} G_{\mu \nu}^a(t,x) G_{\mu \nu}^a(t,x)\,,
  \label{eq:def_E}
\end{align}
are finite after just the renormalization of $g$ and $m_f$.
This allows for a direct comparison of results
obtained in different regularization schemes (lattice and perturbation
theory, for example).

On the contrary, flowed quark fields require a renormalization factor
$Z_\chi^{1/2}(\alphas)$ in order to render Green's functions finite. In
the \msbar\ scheme, it
is\,\cite{Luscher:2010iy,Harlander:2018zpi}
\begin{equation}
  \begin{split}
        Z^{-1}_\chi(\alphas) &= 1
        -\frac{\alphas}{4\pi}\frac{\gamma_{\chi,0}}{2\ep}
        +\left(\frac{\alphas}{4\pi}\right)^2\left[\frac{1}{\ep^2}\left(
          \frac{\gamma_{\chi,0}^2}{8} +
          \frac{\beta_0\gamma_{\chi,0}}{4}\right)
          -\frac{\gamma_{\chi,1}}{4\ep}\right] + \order{\alphas^3}\,,
    %% \label{eq:}
  \end{split}
\end{equation}
with
\begin{equation}
  \begin{split}
  \gamma_{\chi,0} &= 6\,\CF\,,\\
  \gamma_{\chi, 1} &=
  \left( \frac{223}{3} - 16 \ln 2 \right) \CA \CF - \left( 3 + 16
  \ln 2 \right) \CF^2 - \frac{44}{3} \CF \TF\,.
    %% \label{eq:}
  \end{split}
\end{equation}

The quantity
\begin{align}
  S (t,x) \equiv Z_\chi\,\sum_{f=1}^{\NF}\bar{\chi}_f(t,x) \chi_f(t,x)
  \label{eq:def_S}
\end{align}
thus acquires an anomalous dimension, which prevents a direct comparison
of results from different regularization schemes. Alternatively, one may
work with ``ringed quark fields''\,\cite{Makino:2014taa}, which amounts to using
\begin{equation}
  \begin{split}
    \mathring{Z}_\chi(t,\mu) = -\frac{2\NC \NF}{(4\pi t)^2}\, \langle
    R(t)\rangle^{-1}\,,\quad\mbox{with}\quad R(t,x) =
    \sum_{f=1}^{\NF}\bar{\chi}_f(t,x) \overleftrightarrow{\slashed{\mathcal{D}}}^\mathrm{F} \chi_f(t,x)
  \label{eq:Zring}
  \end{split}
\end{equation}
instead of $Z_\chi$ in order to renormalize the quark fields, where
\begin{equation}
  \overleftrightarrow{\mathcal{D}}^\mathrm{F} = \mathcal{D}^\mathrm{F} - \overleftarrow{\mathcal{D}}^\mathrm{F} \,.
\end{equation}
This
corresponds to a ``physical'' renormalization scheme, which means that
the anomalous dimension of the operator
\begin{equation}
  \begin{split}
    \mathring{S}(t,x) = \zeta_\chi(t,\mu)\,S(t,x)\,,\qquad\mbox{with}\quad
    \zeta_\chi(t,\mu) \equiv Z_\chi^{-1}\mathring{Z}_\chi(t,\mu)\,,
    \label{eq:zetachi}
  \end{split}
\end{equation}
vanishes. In \eqn{eq:Zring}, we used the notation
\begin{equation}
  \begin{split}
    \langle \mathcal{O}(t)\rangle \equiv \int\dd^4 x\,
    \langle \mathcal{O}(t,x)\rangle
    %% \label{eq:}
  \end{split}
\end{equation}
for the zero-momentum Fourier transform of the vacuum expectation value
of an operator $\mathcal{O}(t,x)$, which we adopt in the following.

The Feynman rules for the operators $E(t,x)$, $S(t,x)$, and $R(t,x)$
introduced above are obtained in the same manner as those described in
\sct{ssec:Feynman_rules}. Expressing them in terms of the flowed fields
$B(t,x)$ and $\chi_f(t,x)$, the operator $E(t,x)$ results in a bilinear,
trilinear, and a quartic gluon vertex as displayed in
\eqs{fr:egg}--\noeqn{fr:egggg}, $S(t,x)$ corresponds to the single
bilinear quark vertex given in \eqn{fr:sqq}, and $R(t,x)$ to the
bilinear quark vertex of \eqn{fr:rqq} and the quark-gluon vertex of
\noeqn{fr:rqqg}. Similar to the flow vertices, the lines attached to
these vertices are depicted with dashed arrows, indicating that they can
be both flow lines and propagators.  In contrast to the flow vertices,
the flow time of all lines in these vertices is directed towards the
corresponding operator.  Therefore, these vertices set the largest flow
time in the Feynman diagrams.


Using the methods described in \scts{sec:pert} and \ref{sec:implementation}, we can calculate the
vacuum expectation values of $E(t,x)$, $S(t,x)$, and $R(t,x)$ through
three loops. Sample diagrams for these quantities are shown in
\fig{fig:example_diagrams}. The number of diagrams in each case is given
in Table\,\ref{tab:numdias}; recall, however, that this takes into
account the re-writing of the \four-gluon vertex into trilinear vertices
as described in \sct{sec:implementation}. Not included in this number
are diagrams with closed flow-line loops, but integrals with scale-less
sub-loops are counted in.

\begin{figure}
  \begin{center}
    \begin{tabular}{ccc}
      %% \begin{tabular}{c}
         \includegraphics[]{images/E_LO.pdf} (a) &
      %% \end{tabular} &
      %% \begin{tabular}{c}
         \includegraphics[]{images/S_NLO_example.pdf} (b) &
      %% \end{tabular} &
      %% \begin{tabular}{c}
         \includegraphics[]{images/E_NNLO_example.pdf} (c) \\[2em]
      %% \end{tabular} \\
      %% \begin{tabular}{c}
         \includegraphics[]{images/E_NNLO_example2.pdf}\hspace*{-1em}
         (d)\hspace*{1em} &
      %% \end{tabular} &
      %% \begin{tabular}{c}
         \includegraphics[]{images/S_NNLO_example.pdf} (e) &
      %% \end{tabular} &
      %% \begin{tabular}{c}
         \includegraphics[]{images/R_NNLO_example.pdf} (f) \\
      %% \end{tabular} \\
    \end{tabular}
  \end{center}
	\caption{The only diagram contributing to $\langle E(t)\rangle$
          at one loop (a); a \two-loop diagram for $\langle S(t)\rangle$
          and $\langle R(t)\rangle$ (b); and four three-loop diagrams:
          two diagrams for $\langle E(t)\rangle$ (c,d); one diagram for
          $\langle S(t)\rangle$ and $\langle R(t)\rangle$ (e); and one
          diagram for $\langle R(t)\rangle$ (f).}
	\label{fig:example_diagrams}
\end{figure}


\begin{table}
  \begin{center}
    \caption[]{\label{tab:numdias} The number of Feynman diagrams
      contributing to the quantities computed in this section. At
      \three-loop level, the number for ``quenched \qcd'' (marked as
      $\TR=0$) is shown separately.}

    \begin{tabular}{|c||c|c|c|c|}
      \hline
      \#loops & 1 & 2 & 3 & 3  ($\TR=0$)\\\hline\hline
      $\langle E(t)\rangle$ & 1 & 11 & 232 & 211 \\\hline
      $\langle S(t)\rangle$ & 1 & 8 & 210 & 202 \\\hline
      $\langle R(t)\rangle$ & 1 & 11 & 311 & 300 \\\hline
    \end{tabular}
  \end{center}
\end{table}


\subsection{Results for the gluon condensate at three loops}

We write the perturbative result for the gluon condensate as
\begin{equation}
  \langle E(t) \rangle = \frac{3\alphas}{4\pi t^2} \frac{\NA}{8}
  \left[e_0 + \frac{\alphas}{4\pi} e_1 + \left(\frac{\alphas}{4\pi}\right)^2 e_2 +
    \mathcal{O}(\alphas^3)\right] + \mathcal{O}(m)\,,
  \label{eq:et}
\end{equation}
where $\alphas=\alphas(\mu)$, with $\mu$ the renormalization scale. As
indicated, quark mass terms are neglected here.
Their effect at \nlo\ has been studied in \citere{Harlander:2016vzb} and it was found to be negligible.
In this case, the
$e_i$ depend only on the product $z\equiv \mu^2t$. For reasons that will become
clear shortly, it is convenient to parameterize this dependence in
terms of the variable
\begin{equation}
  \begin{split}
    L(z) \equiv \ln(2z) + \EulerGamma\,.
    %% \label{eq:}
  \end{split}
\end{equation}
A
natural choice for the renormalization scale is defined by the inverse
``smearing radius'' $q_8\equiv 1/\sqrt{8t}$ of the gradient
flow. However, from explicit higher order calculations, it was found
that the slightly larger value
\begin{equation}
  \begin{split}
    \mu_0 = \frac{e^{-\EulerGamma/2}}{\sqrt{2t}} \approx 1.5\,q_8\,,
    \label{eq:mu0}
  \end{split}
\end{equation}
corresponding to $L(\mu_0^2t)=0$, improves the stability of the
perturbative corrections\,\cite{Harlander:2018zpi} (see also
\citere{Harlander:2016vzb}). This was corroborated by the behavior of
the $t\to 0$ extrapolation for thermodynamical quantities evaluated in
\citere{Iritani:2018idk} using the \two-loop result of
\citere{Harlander:2018zpi}.

We thus write the coefficients of \eqn{eq:et} as
\begin{equation}
  \begin{split}
    e_0(z) &= e_{0,0}\,,\qquad
    e_1(z) = e_{1,0} + \beta_0\,L(z)\,,\\
    e_2(z) &= e_{2,0} + (2\beta_0\,e_{1,0}+\beta_1)\,L(z) +
    \beta_0^2\,L^2(z)\,,
    \label{eq:efull}
  \end{split}
\end{equation}
with $\beta_0$, $\beta_1$ from \eqn{eq:anombg}.  For the coefficients
$e_{i,j}$, we find
\begin{equation}
  \begin{split}
    e_{0,0} =&\ 1\,,\qquad
    e_{1,0} = \left(\frac{52}{9} + \frac{22}{3} \ln 2 - 3 \ln 3\right)
    \CA - \frac{8}{9} \TF\\
    e_{2,0} =&\ 27.9786\, \CA^2  - (31.5652\ldots)\, \TF\CA
    + \left(16 \zeta(3) - \frac{43}{3}\right)
    \TF\CF + \left(\frac{8\pi^2}{27} -
    \frac{80}{81}\right) \TF^2\,,
    \label{eq:econst}
  \end{split}
\end{equation}
where $\zeta(z)$ is Riemann's $\zeta$ function with $\zeta(3) =
1.20206\ldots$. The three dots in the coefficient of $\TF\CA$ indicate
that we were able to obtain an expression in analytical form. It can be
found in Appendix\,\ref{app:analytic}, \eqn{eq:e20catf}.  Only four
decimal places of the numerical result for the $\CA^2$ coefficient are
displayed here, while our estimate of the numerical accuracy, obtained by
propagating the uncertainty of the individual integrals to the final
result, is about six digits beyond that. This estimate matches the
observed differences between the numerical and analytical results in the
cases where the latter have been computed.  The same statements hold for
the other observables listed below.  The \nlo\ coefficient $e_1$ was
first evaluated in \citere{Luscher:2010iy}.  Setting $\mu=q_8$, one
finds that the \nnlo\ result agrees with \citere{Harlander:2016vzb} at the
sub-percent level.\footnote{To be precise, all color coefficients of
  \citere{Harlander:2016vzb} are compatible with our new calculation,
  except for the $\CA^2$ term, for which ``only'' the first three digits
  agree, corresponding to an overly optimistic uncertainty estimate of
  \citere{Harlander:2016vzb} for that term. The effect for any practical
  application should be irrelevant.}

The logarithmic terms $e_{n,k}L^k(z)$ obey the all-order recursion
formula ($1\leq k\leq n$)
\begin{equation}
  \begin{split}
    k\,e_{n,k} = \sum_{l=0}^{n-1}(n-l)e_{n-l-1,k-1}\beta_l\,,
  \end{split}
\end{equation}
which follows from renormalization-group invariance, i.e.
\begin{equation}
  \begin{split}
    \mu\dderiv{}{\mu}\langle E(t)\rangle = 0\,,\qquad
    \mu^2\dderiv{}{\mu^2}\alphas = \alphas\beta(\alphas) =
    -\frac{\alpha^2_s}{4\pi}\left(\beta_0 + \frac{\alphas}{4\pi}\beta_1
    + \ldots\right)\,.
    \label{eq:rgee}
  \end{split}
\end{equation}
This provides a welcome check of our calculation.  The residual
dependence on $\mu$ can be used as probe of the perturbative behavior.
We refer to \citere{Harlander:2016vzb} for such a study.

Since $\langle E(t) \rangle$ is also gauge independent, we are free to
choose Feynman gauge for the \qcd\ gauge parameter ($\xi=1$) and $\kappa = 1$ for the gradient-flow gauge parameter
for the evaluation of the results
presented above, which facilitates the actual calculation
significantly. However, we also check gauge invariance explicitly by
evaluating the terms with the highest power in the gauge parameter
$\xi$, i.e.\ $\xi^4$, and show that its coefficient vanishes. Indeed,
this already happens algebraically after the reduction to master
integrals, even before their numerical values are inserted.  Keeping the
gauge parameter fully general unfortunately leads to intermediate
expressions which exceed our available computing resources.

\subsection{Results for the quark condensate at three loops}

To obtain the quark condensate with ringed fields we first consider the
conversion factor $\zeta_\chi(t,\mu)$ from \msbar\ renormalized to
ringed quark fields as defined in \eqn{eq:zetachi}. Applying the methods
described in \sct{sec:pert}, we can calculate the diagrams contributing
to the Green's function $\langle R(t)\rangle$ for massless quarks, which
eventually leads to the result
\begin{equation}
	\begin{split}
	  \zeta_\chi(t,\mu) = 1 &+ \frac{\alphas}{4\pi} \left(
          \frac{\gamma_{\chi, 0}}{2} L(\mu^2t) - 3 \CF \ln 3 - 4 \CF \ln
          2 \right) \\ &+ \left(\frac{\alphas}{4\pi}\right)^2 \Bigg\{
          \frac{\gamma_{\chi, 0}}{4} \left( \beta_0 +
          \frac{\gamma_{\chi, 0}}{2} \right) L^2 (\mu^2t) + \Big[
            \frac{\gamma_{\chi, 1}}{2} - \frac{\gamma_{\chi, 0}}{2}
            \left( \beta_0 + \frac{\gamma_{\chi, 0}}{2} \right) \ln 3
            \\ & \qquad \qquad - \frac{2}{3} \gamma_{\chi, 0} \left(
            \beta_0 + \frac{\gamma_{\chi, 0}}{2}\right) \ln 2 \Big]
          L(\mu^2t) + C_2 \Bigg\} + \order{\alphas^3,m} \,.
                \label{eq:zfchi}
	\end{split}
\end{equation}
The \nlo\ result has been obtained in \citere{Makino:2014taa}, and the
$\mu$-dependent \nnlo\ terms were derived in
\citere{Harlander:2018zpi}. That paper also used a preliminary result
for the coefficient $C_2$, derived for the first time in the current
paper. Keeping four decimal places for the individual color coefficients,
we find
\begin{equation}
  C_2 = -23.7947 \, \CA \CF + 30.3914 \, \CF^2 -(3.9226\ldots) \CF
  \TF\,,
  \label{eq:c2}
\end{equation}
where again, as indicated by the three dots, the coefficient of $\CF\TF$
was obtained in analytical form and is given in
Appendix\,\ref{app:analytic}, \eqn{eq:c2cftf}.

Let us now turn to the quark condensate $\langle S(t)\rangle$. Since it
vanishes in the chiral limit $m_f=0$, we compute the leading coefficient
of the expansion in the quark masses $m_f$,
\begin{align}
  \left\langle \zeta_\chi S(t) \right\rangle \equiv \langle
  \mathring{S}(t) \rangle = \sum_{f=1}^{\NF} s_f(t) + \mathcal{O}(m^2)\,,\qquad
  s_f(t) \equiv m_f\dderiv{\langle
    \mathring{S}(t)\rangle}{m_f}\bigg|_{m=0}\,,
  \label{eq:sring}
\end{align}
where it is understood that \textit{all} quark masses are set to zero
after taking the derivative.  This expansion can be applied before the
loop and flow-time integrations are carried out, which means that, in
the quark propagators, we write
\begin{equation}
  \begin{split}
    \frac{1}{\mathrm{i}\slashed{q}+m_f} =
    \frac{1}{\mathrm{i}\slashed{q}} + \frac{m_f}{q^2} +
    \mathcal{O}(m_f^2)\,,
    \label{eq:massexp}
  \end{split}
\end{equation}
and keep only the leading linear term in $m$ of the integrand. The resulting
integrals are then again of the form shown in \eqn{eq:intrep}. We thus
write
\begin{equation}
  s_f(t) = - \frac{\NC m_f}{8\pi^2t} \left[s_0 + \frac{\alphas}{4\pi} s_1 +
    \left(\frac{\alphas}{4\pi}\right)^2 s_2 +
    \mathcal{O}(\alphas^3)\right] + \mathcal{O}(m^2) \,.
\end{equation}
Renormalizing the quark masses in the \msbar\ scheme according to
\eqn{eq:renorm}, we find
\begin{equation}
  \begin{split}
    s_0(z) &= s_{0,0}\,,\qquad
    s_1(z) = s_{1,0} + \frac{\gamma_{m,0}}{2}\,L(z)\,,\\
    s_2(z) &= s_{2,0} + \frac{1}{2}\left[(\gamma_{m,0}+2\beta_0)\,s_{1,0} +
      \gamma_{m,1}\right]L(z) +
    \frac{1}{8}(\gamma_{m,0}+2\beta_0)\gamma_{m,0}\,L^2(z)\,,
    \label{eq:sfull}
  \end{split}
\end{equation}
with $z=\mu^2t$ and the $\gamma_{m,i}$ and $\beta_i$ from \eqn{eq:anombg}.
The non-logarithmic terms are given by
\begin{equation}
  \begin{split}
    s_{0,0} &= 1\,,\qquad s_{1,0} = (4 + 4 \ln 2 - 3 \ln 3) \CF \\
    s_{2,0} &= 19.6422\, \CF^2 + 41.3897\, \CF \CA - (15.7975\dots)\, \TF \CF\,,
    \label{eq:sconst}
  \end{split}
\end{equation}
where again the $\TF\CF$ term is known analytically, as indicated by the
three dots. It can be found in Appendix\,\ref{app:analytic},
\eqn{eq:s20cftf}.  The \nlo\ term has been obtained in
\citere{Makino:2018rys}, the \nnlo\ term is new. Again, only four decimal
places are displayed and our uncertainty estimate for the numerical
integration is several digits beyond that.

The logarithmic terms $s_{n,k}L^k(z)$ obey the all-order recursion formula ($1\leq k\leq n$)
\begin{equation}
  \begin{split}
k\,s_{n,k} = \sum_{l=0}^{n-1}\left[\frac{\gamma_{n,l}}{2} +
  (n-l-1)\beta_l\right]s_{n-l-1,k-1}\,,
  \end{split}
\end{equation}
which follows from renormalization group
invariance, i.e.
\begin{equation}
  \begin{split}
    \mu\dderiv{}{\mu} s_f(t) = 0\,,\qquad
    \mu\dderiv{}{\mu}m_f = m_f\gamma_m(\alphas) =
    - m_f\frac{\alphas}{4\pi}\left(\gamma_{m,0} +
    \frac{\alphas}{4\pi}\gamma_{m,1} + \ldots\right)\,,
    %% \label{eq:}
  \end{split}
\end{equation}
and the renormalization group equation for $\alphas$ given in
\eqn{eq:rgee}. This again provides a welcome check of our calculation.
The residual dependence on $\mu$ can be used as probe of the
perturbative behavior. It is displayed in \fig{fig:S_plot} for three
values of the central energy scale $\mu_0=(2te^{\EulerGamma})^{-1/2}$,
see \eqn{eq:mu0}. We observe a qualitatively similar behavior as was
found for $\langle t^2E(t)\rangle$ in \citere{Harlander:2016vzb}. While
for large $\mu_0$, the reduction of the renormalization scale with
increasing perturbative order is quite significant, this improvement
reduces as $\mu_0$ gets closer to the non-perturbative regime. Taking
the variation of $\mu$ by a factor of two around the central scale as an
estimate of the uncertainty due to missing higher orders, one finds that
they overlap between successive perturbative orders for all three values
of $\mu_0$. However, the trend of the curves does not explicitly
support $\mu_0$ as the central scale in this case, but seems to favor a
slightly smaller scale.

\begin{figure}
  \begin{center}
    \begin{tabular}{cc}
      \includegraphics[width=.45\textwidth]{images/Sring_3.pdf} &
      \includegraphics[width=.45\textwidth]{images/Sring_10.pdf}\\
      \hspace*{2em}(a) & \hspace*{2em}(b)\\
      \multicolumn{2}{c}{\includegraphics[width=.5\textwidth]{%
          images/Sring_130.pdf}}\\
      \multicolumn{2}{c}{\hspace*{2em}(c)}
    \end{tabular}
  \end{center}
	\caption{Renormalization scale dependence of $s_f(t)$, defined
          in \eqn{eq:sring}, for three different values of the central
          scale
          $\mu_0=e^{-\EulerGamma/2}/\sqrt{2t}$. (a)~$\mu_0=3$\,GeV,
          corresponding to $t=(0.03\,\mathrm{fm})^2$, where we use $\NF = 3$
          and $\alpha_\mathrm{s}^{(3)}(3\,\mathrm{GeV}) =
          0.248$;
          (b)~$\mu_0=10$\,GeV ($t=(0.01\,\mathrm{fm})^2$), where $\NF = 5$ and
          $\alpha_\mathrm{s}^{(5)}(10\,\mathrm{GeV}) = 0.178$;
          (c)~$\mu_0=130$\,GeV ($t=(8\cdot
          10^{-4}\,\mathrm{fm})^2$), where $\NF = 5$ and
          $\alpha_\mathrm{s}^{(5)}(130\,\mathrm{GeV}) = 0.112$.
          The
          running of $\alphas(\mu)$ and $m_f(\mu)$ is evaluated with the
          help of
          \texttt{RunDec}\,\cite{Chetyrkin:2000yt,Herren:2017osy} at the
          corresponding loop order. All curves are normalized to the
          \nnlo\ result at $\mu=\mu_0$.}
	\label{fig:S_plot}
\end{figure}

Another important consistency check is that $\langle S(t) \rangle$ and
$\langle R(t) \rangle$ can be renormalized with the same flowed-quark
field renormalization constant $Z_\chi$ as obtained in
\citere{Harlander:2018zpi}.  $\langle S(t) \rangle$ and $\langle
R(t)\rangle$ are also gauge independent.  For the calculation of the
full \nnlo\ result, we again chose Feynman gauge ($\xi=1$) and
$\kappa=1$.  We check gauge invariance explicitly by evaluating the
terms with the highest power in the gauge parameter $\xi$, which is
$\xi^3$ for $\langle S(t) \rangle$ and $\langle R(t)\rangle$, and show
that their coefficients vanish.  This happens here also immediately
after the reduction to master integrals.
